
This study set out to verify that semantic entropy achieves higher AUROC than token-probability for correctness prediction in Llama-3-8B-Base on TriviaQA and NaturalQuestions, as a prerequisite for a scale-dependent interaction hypothesis. The verification failed: SE achieves AUROC of 0.4735 on TriviaQA and 0.5524 on NaturalQuestions, compared to 0.6835 and 0.6551 for token-probability on the same datasets. KLE performs at 0.2642 and 0.3753 respectively, and SelfCheckGPT-BERTScore yields a constant AUROC of 0.5000.

The degenerate_fraction diagnostic records that 89.4\% of TriviaQA queries and 84.8\% of NaturalQuestions queries produce K=1 semantic clusters under N=10 sampling at temperature=1.0 from Llama-3-8B-Base. This observation is consistent with the SE failure: a method that relies on semantic diversity across samples cannot produce discriminative uncertainty scores when the vast majority of queries yield semantically identical samples.

The contribution of this study is primarily diagnostic. The \texttt{degenerate_fraction} metric makes explicit a precondition that is implicit in clustering-based UQ methods and was not reported in their original evaluations. Token-probability is shown to provide a valid uncertainty signal in the base-model, factual-QA regime studied here. The finding that SE superiority over token-probability is not observed for Llama-3-8B-Base under standard evaluation conditions documents a boundary on the known validity of this method.

Several important questions remain open: whether SE functions correctly on instruction-tuned variants of Llama-3; whether temperature adjustment can reduce degenerate_fraction sufficiently to restore SE discriminability on base models; and whether 70B scale results differ from 8B. These are identified as the immediate priorities for follow-up.

